% These computations have not been included in the main paper.

\section{Equivariance of the $G$-Convolution}\label{sec:g-conv}


The $G$-convolution is \textit{equivariant} to the action of the group $G$ on the domain of the signal $\signal$~\cite{cohen2016group}. This means that any action $\action_h$ on the domain of $\signal$ results in a corresponding action $\action'_h$ on the output of the convolutional layer. Specifically, for a filter $\filter $, the equivariance property is:
\begin{equation}
    \filter * \action_{h}[\signal] = \action'_{h}[\filter * \signal], \qquad \text{for all ${h} \in G$,}
\end{equation}
where $\action_h$ and $\action'_h$ represent the actions of the same group element $h$ on $\signal$ and $\filter * \signal$ respectively.



\begin{proof}
Consider a signal $\signal: \domain \rightarrow \mathbb{R}^c$ and a filter $\filter: \domain \rightarrow \mathbb{R}^c$ defined on a domain $\domain$ and codomain $\mathbb{R}^c$, with $c$ the number of channels.

For transformations $h, g \in G$, we have:
\begin{align*}
    \filter * \action_{h}[\signal](g)
        &= \int_{u \in \domain}
            \filter(\action_{g^{-1}}(u)) \action_{h}[\signal](u) du \\
        &= \int_{u \in \domain}
            \filter(\action_{g^{-1}}(u))\signal( \action_{h^{-1}}(u)) du \\
        &= \int_{u' \in \domain}
            \filter(\action_{g^{-1}}(\action_{h}(u')))\signal(u') du'
            \qquad
            \text{($u' =  \action_{h^{-1}}(u)$, i.e., $u = \action_{h}(u')$)} \\
        &= \int_{u \in \domain}
            \filter(\action_{g^{-1}h}(u))\signal(u) du \\
        &= \int_{u \in \domain}
            \filter(\action_{(h^{-1}g)^{-1}}(u))\signal(u) du \\
        &= \int_{u \in \domain}
            \action_{(h^{-1}g)}[\filter](u)\signal(u) du \\
        &=\action'_{h}[\filter * \signal](g)
\end{align*}
where the action $L'$ uses $h$ to act on the domain $G$ of the output signal  $[\filter * \signal]$ by left multiplication with $h^{-1}$.
\end{proof}

This property applies for the $G$-convolutions of the first layer and of the next layers.

\section{Example of Triple Correlation on $C_4$}

In this appendix, we work out the computations of the TC for the group $C_4$ shown in the explanatory Figure of the main text.

The group is $G = C_4$ and the first input real signal is $\Theta$ defined as:

\begin{equation}
    \Theta(0) = 0, \quad \Theta(1) = -1, \quad \Theta(2) = 1, \quad \Theta(3) = 2.
\end{equation}

We compute the elements of the triple correlation.

First the $(0, 0)$ diagonal element.
\scriptsize
\begin{align*}
    \tau_\Theta(0, 0)
    = \sum_{g \in G} \Theta(g)^3
    = 0^3 + (-1)^3 + 1^3 + 2^3
    = 8
\end{align*}

\normalsize
Then we populate the first line of the TC and add the values of the TC for the symmetric elements.
\scriptsize
\begin{align*}
\tau_\Theta(0, 1)
    &= \sum_{g \in G} \Theta(g)^2 \Theta(g+1)\\
    &= \Theta(0)^2 \Theta(0+1)
    + \Theta(1)^2 \Theta(1+1)
    + \Theta(2)^2 \Theta(2+1)
    + \Theta(3)^2 \Theta(3+1) \\
    &= \Theta(0)^2 \Theta(1)
    + \Theta(1)^2 \Theta(2)
    + \Theta(2)^2 \Theta(3)
    + \Theta(3)^2 \Theta(0) \\
    &= 0 + (-1)^2.1 + 1^2.2+0\\
    &=3 \\
    &=\tau_\Theta(1, 0)\\
    &=\tau_\Theta(0^{-1}, 1.0^{-1}) = \tau_\Theta(0, 1) \\
    &=\tau_\Theta(1^{-1}, 0.1^{-1}) = \tau_\Theta(3, 3)\\
\tau_\Theta(0, 2)
    &= \sum_{g \in G} \Theta(g)^2 \Theta(g+2)\\
    &= \Theta(0)^2 \Theta(0+2)
    + \Theta(1)^2 \Theta(1+2)
    + \Theta(2)^2 \Theta(2+2)
    + \Theta(3)^2 \Theta(3+2) \\
    &= \Theta(0)^2 \Theta(2)
    + \Theta(1)^2 \Theta(3)
    + \Theta(2)^2 \Theta(0)
    + \Theta(3)^2 \Theta(1) \\
    &= 0 + (-1)^2.2 + 0 + 2^2.(-1)\\
    &= -2 \\
    &=\tau_\Theta(2, 0)\\
    &=\tau_\Theta(0^{-1}, 2.0^{-1}) = \tau_\Theta(0, 2) \\
    &=\tau_\Theta(2^{-1}, 0.2^{-1}) = \tau_\Theta(2, 2)\\
\tau_\Theta(0, 3)
    &= \sum_{g \in G} \Theta(g)^2 \Theta(g+3)\\
    &= \Theta(0)^2 \Theta(0+3)
    + \Theta(1)^2 \Theta(1+3)
    + \Theta(2)^2 \Theta(2+3)
    + \Theta(3)^2 \Theta(3+3) \\
    &= \Theta(0)^2 \Theta(3)
    + \Theta(1)^2 \Theta(0)
    + \Theta(2)^2 \Theta(1)
    + \Theta(3)^2 \Theta(2) \\
    &= 0 + 0 + 1^2.(-1) + 2^2.1\\
    &= 3 \\
    &=\tau_\Theta(3, 0)\\
    &=\tau_\Theta(0^{-1}, 3.0^{-1}) = \tau_\Theta(0, 3) \\
    &=\tau_\Theta(3^{-1}, 0.3^{-1}) = \tau_\Theta(1, 1).
\end{align*}

\normalsize
Then, we populate the second line of the TC.
\scriptsize
\begin{align*}
\tau_\Theta(1, 2)
    &= \sum_{g \in G} \Theta(g) \Theta(g+1)\Theta(g+2)\\
    &= \Theta(0) \Theta(0+1)\Theta(0+2)
    + \Theta(1) \Theta(1+1)\Theta(1+2)
    + \Theta(2) \Theta(2+1)\Theta(2+2)
    + \Theta(3) \Theta(3+1)\Theta(3+2) \\
    &= \Theta(0) \Theta(1)\Theta(2)
    + \Theta(1) \Theta(2)\Theta(3)
    + \Theta(2) \Theta(3)\Theta(0)
    + \Theta(3) \Theta(0)\Theta(1) \\
    &= 0
    + -1.1.2
    + 0
    + 0\\
    &=-2\\
    &=\tau_\Theta(2, 1)\\
    &=\tau_\Theta(1^{-1}, 2+1^{-1}) = \tau_\Theta(3, 2+3)=\tau_\Theta(3, 1) =\tau_\Theta(1, 3) \\
    &=\tau_\Theta(2^{-1}, 1+2^{-1}) = \tau_\Theta(2, 1+2)=\tau_\Theta(2, 3) =\tau_\Theta(3, 2).
\end{align*}
\normalsize
\section{Example of Triple Correlation on $\mathbb{Z}_2 \times \mathbb{Z}_2$}

In this appendix, we work out the computations of the TC for the group $\mathbb{Z}_2 \times \mathbb{Z}_2$ shown in the explanatory Figure of the main text.

The group is $G = \mathbb{Z}_2 \times \mathbb{Z}_2 $ and the first input real signal is $\Theta$ defined as:
\begin{equation}
    \Theta(00) = 2, \quad \Theta(01) = 0, \quad \Theta(10) = 1, \quad \Theta(11) = 1.
\end{equation}

We compute the elements of the triple correlation.

\scriptsize
\begin{align*}
    \tau_\Theta(00, 00)
    &= \sum_{g \in G} \Theta(g)^3
    = 8+0+1+1
    = 10 \\
    \tau_\Theta(00, 01)
    &=\sum_{g \in G} \Theta(g)^2\Theta(g \cdot 01) \\
    &= \Theta(00)^2\Theta(01)
        + \Theta(01)^2\Theta(01\cdot 01)
        + \Theta(10)^2\Theta(10\cdot 01)
        + \Theta(11)^2\Theta(11\cdot 01)\\
    &= \Theta(00)^2\Theta(01)
        + \Theta(01)^2\Theta(00)
        + \Theta(10)^2\Theta(11)
        + \Theta(11)^2\Theta(10)\\
    &= 2^2 .0 +0^2.2+1^2.1+1^2.1
    = 2 \\
    \tau_\Theta(00, 10)
    &=\sum_{g \in G} \Theta(g)^2\Theta(g \cdot 10) \\
    &= \Theta(00)^2\Theta(10)
        + \Theta(01)^2\Theta(01\cdot 10)
        + \Theta(10)^2\Theta(10\cdot 10)
        + \Theta(11)^2\Theta(11\cdot 10)\\
    &= \Theta(00)^2\Theta(10)
        + \Theta(01)^2\Theta(11)
        + \Theta(10)^2\Theta(00)
        + \Theta(11)^2\Theta(01)\\
    &= 2^2 .1 +0^2.1+1^2.2+1^2.0
    = 3 \\
    \tau_\Theta(00, 11)
    &= \sum_{g \in G} \Theta(g)^2\Theta(g \cdot 11) \\
    &= \Theta(00)^2\Theta(11)
        + \Theta(01)^2\Theta(01\cdot 11)
        + \Theta(10)^2\Theta(10\cdot 11)
        + \Theta(11)^2\Theta(11\cdot 11)\\
    &= \Theta(00)^2\Theta(11)
        + \Theta(01)^2\Theta(10)
        + \Theta(10)^2\Theta(01)
        + \Theta(11)^2\Theta(00)\\
    &= 2^2 .1 +0^2.1+1^2.0+1^2.2 = 6 \\
    % Don't compute these: Get the diagonal elements with symmetry
    % \tau_\Theta(01, 01)
    % &= \sum_{g \in G} \Theta(g)\Theta(g \cdot 01)^2 \\
    % &= \Theta(00)\Theta(00 \cdot 01)^2
    % + \Theta(01)\Theta(01 \cdot 01)^2
    % + \Theta(10)\Theta(10 \cdot 01)^2
    % + \Theta(11)\Theta(11 \cdot 01)^2 \\
    % &= \Theta(00)\Theta(01)^2
    % + \Theta(01)\Theta(00)^2
    % + \Theta(10)\Theta(11)^2
    % + \Theta(11)\Theta(10)^2 \\
    % &= 2.0^2
    % + 0.2^2
    % + 1.1^2
    % + 1.1^2 = 2 \\
    % \tau_\Theta(10, 10)
    % &= \sum_{g \in G} \Theta(g)\Theta(g \cdot 10)^2 \\
    % &= \Theta(00)\Theta(00 \cdot 10)^2
    % + \Theta(01)\Theta(01 \cdot 10)^2
    % + \Theta(10)\Theta(10 \cdot 10)^2
    % + \Theta(11)\Theta(11 \cdot 10)^2 \\
    % &= \Theta(00)\Theta(10)^2
    % + \Theta(01)\Theta(11)^2
    % + \Theta(10)\Theta(00)^2
    % + \Theta(11)\Theta(01)^2 \\
    % &= 2.1^2
    % + 0.1^2
    % + 1.2^2
    % + 1.0^2 = 6 \\
    % \tau_\Theta(11, 11)
    % &= \sum_{g \in G} \Theta(g)\Theta(g \cdot 11)^2 \\
    % &= \Theta(00)\Theta(00 \cdot 11)^2
    % + \Theta(01)\Theta(01 \cdot 11)^2
    % + \Theta(10)\Theta(10 \cdot 11)^2
    % + \Theta(11)\Theta(11 \cdot 11)^2 \\
    % &= \Theta(00)\Theta(11)^2
    % + \Theta(01)\Theta(10)^2
    % + \Theta(10)\Theta(01)^2
    % + \Theta(11)\Theta(00)^2 \\
    % &= 2.1^2
    % + 0.1^2
    % + 1.0^2
    % + 1.2^2 = 6 \\
    \tau_\Theta(01, 10)
    &= \sum_{g \in G} \Theta(g)\Theta(g \cdot 01) \Theta(g \cdot 10) \\
    &= \Theta(00)\Theta(00 \cdot 01) \Theta(00 \cdot 10)
    + \Theta(01)\Theta(01 \cdot 01) \Theta(01 \cdot 10)
    + \Theta(10)\Theta(10 \cdot 01) \Theta(10 \cdot 10)
    + \Theta(11)\Theta(11 \cdot 01) \Theta(11 \cdot 10)\\
    &=\Theta(00)\Theta( 01) \Theta(10)
    + \Theta(01)\Theta(00) \Theta(11)
    + \Theta(10)\Theta(11) \Theta(00)
    + \Theta(11)\Theta(10) \Theta(01)\\
    &=2.0.1
    + 0.2.1
    + 1.1.2
    + 1.1.0\\
    &=2 \\
    \tau_\Theta(11, 10)
    &= \sum_{g \in G} \Theta(g)\Theta(g \cdot 11) \Theta(g \cdot 10) \\
    &= \Theta(00)\Theta(00 \cdot 11) \Theta(00 \cdot 10)
    + \Theta(01)\Theta(01 \cdot 11) \Theta(01 \cdot 10)
    + \Theta(10)\Theta(10 \cdot 11) \Theta(10 \cdot 10)
    + \Theta(11)\Theta(11 \cdot 11) \Theta(11 \cdot 10)\\
    &=\Theta(00)\Theta(11) \Theta(10)
    + \Theta(01)\Theta(10) \Theta(11)
    + \Theta(10)\Theta(01) \Theta(00)
    + \Theta(11)\Theta(00) \Theta(01)\\
    &=2.1.1
    + 0.1.1
    + 1.0.2
    + 1.2.0\\
    &=2
\end{align*}

\normalsize

The second input real signal is $\Theta'$ defined as:
\begin{equation}
    \Theta'(00) = 0, \quad \Theta'(01) = 1, \quad \Theta'(10) = 1, \quad \Theta'(11) = 2.
\end{equation}

We compute the triple correlation for $\Theta'$

\scriptsize
\begin{align*}
    \tau_\Theta'(00, 00)
    &= \sum_{g \in G} \Theta'(g)^3
    = 0+1+1+8
    = 10 \\
    \tau_\Theta'(00, 01)
    &=\sum_{g \in G} \Theta'(g)^2\Theta'(g \cdot 01) \\
    &= \Theta'(00)^2\Theta'(01)
        + \Theta'(01)^2\Theta'(01\cdot 01)
        + \Theta'(10)^2\Theta'(10\cdot 01)
        + \Theta'(11)^2\Theta'(11\cdot 01)\\
    &= \Theta'(00)^2\Theta'(01)
        + \Theta'(01)^2\Theta'(00)
        + \Theta'(10)^2\Theta'(11)
        + \Theta'(11)^2\Theta'(10)\\
    &= 0^2 .1 +1^2.0+1^2.2+2^2.1
    = 6 \\
    \tau_\Theta'(00, 10)
    &=\sum_{g \in G} \Theta'(g)^2\Theta'(g \cdot 10) \\
    &= \Theta'(00)^2\Theta'(10)
        + \Theta'(01)^2\Theta'(01\cdot 10)
        + \Theta'(10)^2\Theta'(10\cdot 10)
        + \Theta'(11)^2\Theta'(11\cdot 10)\\
    &= \Theta'(00)^2\Theta'(10)
        + \Theta'(01)^2\Theta'(11)
        + \Theta'(10)^2\Theta'(00)
        + \Theta'(11)^2\Theta'(01)\\
    &= 2^2 .1 +1^2.2+1^2.0+2^2.1
    = 14 \\
    \tau_\Theta'(00, 11)
    &= \sum_{g \in G} \Theta'(g)^2\Theta'(g \cdot 11) \\
    &= \Theta'(00)^2\Theta'(11)
        + \Theta'(01)^2\Theta'(01\cdot 11)
        + \Theta'(10)^2\Theta'(10\cdot 11)
        + \Theta'(11)^2\Theta'(11\cdot 11)\\
    &= \Theta'(00)^2\Theta'(11)
        + \Theta'(01)^2\Theta'(10)
        + \Theta'(10)^2\Theta'(01)
        + \Theta'(11)^2\Theta'(00)\\
    &= 0^2 .2 +1^2.1+1^2.1+2^2.0 = 2 \\
    \tau_\Theta'(01, 10)
    &= \sum_{g \in G} \Theta'(g)\Theta'(g \cdot 01) \Theta'(g \cdot 10) \\
    &= \Theta'(00)\Theta'(00 \cdot 01) \Theta'(00 \cdot 10)
    + \Theta'(01)\Theta'(01 \cdot 01) \Theta'(01 \cdot 10)
    + \Theta'(10)\Theta'(10 \cdot 01) \Theta'(10 \cdot 10)
    + \Theta'(11)\Theta'(11 \cdot 01) \Theta'(11 \cdot 10)\\
    &=\Theta'(00)\Theta'( 01) \Theta'(10)
    + \Theta'(01)\Theta'(00) \Theta'(11)
    + \Theta'(10)\Theta'(11) \Theta'(00)
    + \Theta'(11)\Theta'(10) \Theta'(01)\\
    &=0.1.1
    + 1.0.1
    + 1.2.0
    + 2.1.1\\
    &=2 \\
    \tau_\Theta'(11, 10)
    &= \sum_{g \in G} \Theta'(g)\Theta'(g \cdot 11) \Theta'(g \cdot 10) \\
    &= \Theta'(00)\Theta'(00 \cdot 11) \Theta'(00 \cdot 10)
    + \Theta'(01)\Theta'(01 \cdot 11) \Theta'(01 \cdot 10)
    + \Theta'(10)\Theta'(10 \cdot 11) \Theta'(10 \cdot 10)
    + \Theta'(11)\Theta'(11 \cdot 11) \Theta'(11 \cdot 10)\\
    &=\Theta'(00)\Theta'(11) \Theta'(10)
    + \Theta'(01)\Theta'(10) \Theta'(11)
    + \Theta'(10)\Theta'(01) \Theta'(00)
    + \Theta'(11)\Theta'(00) \Theta'(01)\\
    &=0.2.1
    + 1.1.2
    + 1.1.0
    + 1.0.1\\
    &=2
\end{align*}
\normalsize
\newpage

\section{Algorithmic Reduction of the Triple Correlation}\label{app:reduction}

\subsection{Direct Approach}

In the ``direct approach'', we ask:\\

\textit{which TC coefficients do we need to reconstruct the original signal?}\\

We note that we do not care about whether the reconstruction is computationally expensive, as we will not implement it in practice. We are only interested in knowing the minimal number of TC coefficients that fully represent $f$.

We consider a simple example on $Z_n$

Consider a real signal defined on $Z_n$: $f(0), ..., f(n-1)$. Try to design an algorithm to recover that signal from some values of the triple correlation.


We compute some quantities related to the TC that could be helpful:
\begin{equation}
    a(0, 0) = \sum_{x = 0}^{n-1} f(x)^3
\end{equation}


\begin{equation}
   \sum_{x'=0}^{n-1} a(0, x') =  \sum_{x'=0}^{n-1} \sum_{x = 0}^{n-1} f(x)^2f(x-x')
\end{equation}


\subsection{Bispectrum Approach}

In the ``bispectrum approach'', we ask:\\

\textit{which TC coefficients do we need to reconstruct the first column of the bispectrum?}\\

If we can reconstruct the first column of the bispectrum, then we can reconstruct the signal. Again, we do not care whether the reconstruction itself is computationally expensive.

We write the definition of the bispectrum (first column) and look at the triple correlation coefficients involved.

We have:
\begin{align*}
    \beta(f)_{k_1, k_2}
        &=
    \frac{1}{n^3}\sum_{x_1 \in Z_n} \sum_{x_2 \in Z_n}
    T(f)(x_1, x_2)\left[
                \rho_1(x_1)^{\dagger} \otimes \rho_2(x_2)^{\dagger}
            \right]\\
    &=
    \frac{1}{n^3}\sum_{x_1 \in Z_n} \sum_{x_2 \in Z_n}
    \left(\sum_{x \in Z_n}
       f^*(x)
       f\left(x + x_1\right) f\left(x+ x_2\right)\right)
            \left[
                \rho_1(x_1)^{\dagger} \otimes \rho_2(x_2)^{\dagger}
            \right]
\end{align*}
where each representation $\rho_j$ is associated with and integer $k_j$ such that: $\rho_1(x') = \exp(-i\frac{2k_1\pi x'}{n})$ and $\rho_2(x') = \exp(-i\frac{2k_2\pi x'}{n})$ for any $x' \in Z_n$.

We are only interested in the first column of the bispectrum. Thus, we set $k_1 = 0$ or equivalently $\rho_1(x') = 1$ for all $x' \in Z_n$:
\begin{align*}
    \beta(f)_{0, k_2}
        &=
    \sum_{x_1 \in Z_n} \sum_{x_2 \in Z_n}
    T(f)(x_1, x_2)\left[
                1 \otimes \rho_2(x_2)^{\dagger}
            \right]\\
    &=
    \sum_{x_1 \in Z_n} \sum_{x_2 \in Z_n}
    \left(\sum_{x \in Z_n}
       f^*(x)
       f\left(x + x_1\right) f\left(x+ x_2\right)\right)
            \left[
                1\otimes \rho_2(x_2)^{\dagger}
            \right]\\
    % &=
    % \sum_{x_1 \in Z_n}
    % \left(\sum_{x \in Z_n}
    %    f^*(x)
    %    f\left(x + x_1\right) \sum_{x_2 \in Z_n} f\left(x+ x_2\right) \rho_2(x_2)^{\dagger}\right).
\end{align*}

At this stage, we see that even if we set $k_1=0$, we still need all the $x_1, x_2$ terms of the triple correlation.

We can continue the computations, to see if we can get rid of $x_1$.
\begin{align*}
    \beta(f)_{0, k_2}
    &=
   \sum_{x_2 \in Z_n} \left(\sum_{x \in Z_n} f^*(x) f\left(x+ x_2\right) \rho_2(x_2)^{\dagger}
      \sum_{x_1 \in Z_n}  f\left(x + x_1\right) \right)\\
          &=
    \sum_{x_2 \in Z_n} \left(\sum_{x \in Z_n}  f^*(x) f\left(x+ x_2\right) \rho_2(x_2)^{\dagger}
      \sum_{x_1' \in Z_n}  f\left(x_1'\right) \right) \quad \text{(change of variable $x'_1 = x + x_1$)}\\
          &=
    \sum_{x_1' \in Z_n}\sum_{x_2 \in Z_n}\left(\sum_{x \in Z_n}  f^*(x) f\left(x+ x_2\right) f\left(x_1'\right)\right) \rho_2(x_2)^{\dagger},
\end{align*}

Unfortunately, this does not make a simpler TC term appear.

We try another version of the computations, to see if we can get rid of $x_2$.
\begin{align*}
    \beta(f)_{0, k_2}
    &=
    \sum_{x_1 \in Z_n}
    \left(\sum_{x \in Z_n}
       f^*(x)
       f\left(x + x_1\right) \rho_2(-x).\hat{f}_{k_2}\right) \quad \text{(equivariance of Fourier transform)}\\
    &=
   \hat{f}_{k_2}
    \left(\sum_{x \in Z_n}
       f^*(x)  \rho_2(-x)\sum_{x_1 \in Z_n}
       f\left(x + x_1\right) \right) \\
    &=
    \hat{f}_{k_2}
    \left(\sum_{x \in Z_n}
       f^*(x)  \rho_2(-x)\sum_{x_1' \in Z_n}
       f\left(x_1'\right) \right) \quad \text{change of variable}\\
    &=
    \hat{f}_{k_2}
    \left(\sum_{x \in Z_n}
       f(x)  \rho_2(x)^\dagger\right)^*\left(\sum_{x_1' \in Z_n}
       f\left(x_1'\right) \right)\\
    &=
    \hat{f}_{k_2}
    \hat{f}_{k_2}^*\left(\sum_{x_1' \in Z_n}
       f\left(x_1'\right) \right)\\
    &=
    \hat{f}_{k_2}
    \hat{f}_{k_2}^*\hat{f}_{0}
\end{align*}
We recovered the formula of the bispectrum, but this does not help figuring which TC coefficients we need: we find that we need only some fourier coefficients, as expected.

\paragraph{From previous meeting} Define the triple correlation of a single group element:

\begin{align}
\label{eq:tc0}
    a(0, g) = \int_{g^{\prime}} f(g^{\prime})^2 f(g^{\prime} g) dg^{\prime} \\
\end{align}

Recall the bispectrum on a single irrep:

\begin{align}
    A(0, \rho) &= \hat{f}_0 \hat{f}_{\rho} \hat{f}_{\rho}^{\dagger} \\
    &= \int_{g^{\prime}} f(g^{\prime}) dg^{\prime} \int_g f(g) \rho(g)^{\dagger} dg \int_g f(g)^{\dagger} \rho(g) dg
\end{align}

We take the group Fourier transform of \ref{eq:tc0} to obtain:

\begin{align}
    \mathcal{F}_{\rho}(a) &= \int_{g} a(0, g) \rho_{g}^{\dagger} dg \\
    &= \int_{g} \int_{g_2} f(g_2)^2 f(g_2 g) dg_2 \rho_{g}^{\dagger} dg \\
    &= \int_{g} \int_{g_2} f(g_2) \rho_g f(g_2) \rho_g^{\dagger} f(g_2 g) dg_2 \rho_{g}^{\dagger} dg \\
    &= \int_{g_2} f(g_2) f(g_2) dg_2  \int_{g} f(g_2 g)\rho_g  \rho_g^{\dagger}\rho_{g}^{\dagger} dg \\
    &= \int_{g_2} f(g_2) f(g_2) dg_2  \int_{g'} f(g')\rho_{g_2}\rho_{g'} dg' \\
&= \int_{g_2} f(g_2) f(g_2) dg_2  \rho_{g_2} \left(\int_{g'} f(g')\rho_{g'} dg'\right) \\
&= \int_{g'} f(g')f(g')  \rho_{g'} dg'  \left(\int_{g} f(g)\rho_{g} dg\right)
\end{align}
if $\rho_{g'} = \rho$ was independent of $g'$, we could conclude, because we could say that the function $f$ and $\rho f$ are linearly dependant, such that they are colinear vectors in the space of functions. We could then apply the equality case of Cauchy-Schwarz to get:
\begin{align}
\mathcal{F}_{\rho}(a)
    &= \int_{g'} f(g')f(g')  \rho dg'  \left(\int_{g} f(g)\rho_{g} dg\right) \\
    &= \sqrt{\int_{g'} f(g')^2dg'} \sqrt{\int_{g'} \rho^2 f(g')^2dg'}  \left(\int_{g} f(g)\rho_{g} dg\right)
\end{align}
which is close to what we want, but not even what we want.


% \begin{proof}
% Consider a signal $\signal$ and a filter $\phi_k$. For a transformation ${\tilde{g}} \in G$, we have:
% \begin{align*}
%     \filter * \action_{\tilde{g}}[\signal](g)
%         &= \int_{u \in \domain}
%             \filter(\action_{g^{-1}}(u)) \action_{\tilde{g}}[\signal](u) du \\
%         &= \int_{u \in \domain}
%             \filter(\action_{g^{-1}}(u))\signal( \action_{\tilde{g}}(u)) du \\
%         &= \int_{u' \in \domain}
%             \filter(\action_{g^{-1}}(\action_{\tilde{g}^{-1}}(u')))\signal(u') du'
%             \qquad
%             \text{($u' =  \action_{\tilde{g}}(u)$ i.e. $u =  \action_{\tilde{g}^{-1}}(u)$)} \\
%         &= \int_{u \in \domain}
%             \filter(\action_{g^{-1}\tilde{g}^{-1}}(u))\signal(u) du \\
%         &= \int_{u \in \domain}
%             \filter(\action_{(\tilde{g}g)^{-1}}(u))\signal(u) du \\
%         &=\action'_{\tilde{g}}[\filter * \signal](g)
% \end{align*}
% where the action $L'$ uses $\tilde{g}$ to act on the domain $G$ of the output signal  $[\filter * \signal]$ by left multiplication.

% TODO: What property allows to say that $g \domain = \domain$? Homogeneity implies that fact to be true, but we can find a more relaxed property and define the convolution on more general spaces.
% \end{proof}
